\documentclass[UTF8]{ctexart}
\usepackage{geometry}
\usepackage{amsmath,amssymb,amsfonts}
\usepackage{array,booktabs}
\geometry{a4paper,margin=2.5cm}

\begin{document}

\begin{center}
\LARGE\bfseries 第三章作业题单
\end{center}

\section*{教材《概率论与数理统计》}

\subsection*{第三章习题（pp.84--88）}

\begin{enumerate}
\renewcommand{\labelenumi}{\textbf{\arabic{enumi}.}}
\item 在一箱子中装有12只开关，其中2只是次品，在其中取两次，每次任取一只，考虑两种试验：(1) 放回抽样；(2) 不放回抽样。我们定义随机变量$X,Y$如下：
\[
X=\begin{cases}
0, & \text{若第一次取出的是正品},\\[2pt]
1, & \text{若第一次取出的是次品};
\end{cases}
\qquad
Y=\begin{cases}
0, & \text{若第二次取出的是正品},\\[2pt]
1, & \text{若第二次取出的是次品}.
\end{cases}
\]
试分别就(1)、(2)两种情况，写出$X$和$Y$的联合分布律。

\setcounter{enumi}{4}
\item 设随机变量$(X,Y)$具有分布函数
\[
F(x,y)=
\begin{cases}
1-\mathrm{e}^{-x}-\mathrm{e}^{-y}+\mathrm{e}^{-x-y}, & x>0,\ y>0,\\[4pt]
0, & \text{其他}.
\end{cases}
\]
求边缘分布函数。

\item 将一枚硬币掷3次，以$X$表示前2次中出现$H$的次数，以$Y$表示3次中出现$H$的次数。求$X,Y$的联合分布律以及$(X,Y)$的边缘分布律。

\setcounter{enumi}{8}
\item 设二维随机变量$(X,Y)$的概率密度为
\[
f(x,y)=
\begin{cases}
cx^{2}y, & x^{2}\le y\le 1,\\[4pt]
0, & \text{其他}.
\end{cases}
\]
\begin{enumerate}
\item[(1)] 确定常数$c$。
\item[(2)] 求边缘概率密度。
\end{enumerate}

\setcounter{enumi}{10}
\item 以$X$记某医院一天出生的婴儿的个数，$Y$记其中男婴的个数，设$X$和$Y$的联合分布律为
\[
P\{X=n,Y=m\}=\frac{\mathrm{e}^{-14}\times 7.14^{m}\times 6.86^{\,n-m}}{m!\,(n-m)!},
\qquad m=0,1,2,\cdots,n;\quad n=0,1,2,\cdots.
\]
\begin{enumerate}
\item[(1)] 求边缘分布律。
\item[(2)] 求条件分布律。
\item[(3)] 特别，写出当$X=20$时，$Y$的条件分布律。
\end{enumerate}

\setcounter{enumi}{12}
\item 在第9题中：
\begin{enumerate}
\item[(1)] 求条件概率密度$f_{X\mid Y}(x\mid y)$，特别，写出当$Y=\dfrac{1}{2}$时$X$的条件概率密度。
\item[(2)] 求条件概率密度$f_{Y\mid X}(y\mid x)$，特别，分别写出当$X=\dfrac{1}{3},X=\dfrac{1}{2}$时$Y$的条件概率密度。
\item[(3)] 求条件概率
\[
P\left\{Y\ge\frac{1}{4}\,\middle|\,X=\frac{1}{2}\right\},\qquad
P\left\{Y\ge\frac{3}{4}\,\middle|\,X=\frac{1}{2}\right\}.
\]
\end{enumerate}

\setcounter{enumi}{14}
\item 设随机变量$X\sim U(0,1)$，当给定$X=x$时，随机变量$Y$的条件概率密度为
\[
f_{Y\mid X}(y\mid x)=
\begin{cases}
x, & 0<y<\dfrac{1}{x},\\[6pt]
0, & \text{其他}.
\end{cases}
\]
\begin{enumerate}
\item[(1)] 求$X$和$Y$的联合概率密度$f(x,y)$。
\item[(2)] 求边缘概率密度$f_{Y}(y)$，并画出它的图形。
\item[(3)] 求$P\{X>Y\}$。
\end{enumerate}

\setcounter{enumi}{16}
\item 
\begin{enumerate}
\item[(1)] 设随机变量$(X,Y)$具有分布函数
\[
F(x,y)=
\begin{cases}
(1-\mathrm{e}^{-ax})\,y, & x\ge 0,\ 0\le y\le 1,\\[4pt]
1-\mathrm{e}^{-ax}, & x\ge 0,\ y>1,\\[4pt]
0, & \text{其他},
\end{cases}
\qquad a>0.
\]
证明$X,Y$相互独立。
\item[(2)] 设随机变量$(X,Y)$具有分布律
\[
P\{X=x,Y=y\}=p^{2}(1-p)^{x+y-2},\quad 0<p<1,\ x,\ y\ \text{均为正整数},
\]
问$X,Y$是否相互独立？
\end{enumerate}

\item 设$X$和$Y$是两个相互独立的随机变量，$X$在区间$(0,1)$上服从均匀分布，$Y$的概率密度为
\[
f_{Y}(y)=
\begin{cases}
\dfrac{1}{2}\mathrm{e}^{-y/2}, & y>0,\\[6pt]
0, & y\le 0.
\end{cases}
\]
\begin{enumerate}
\item[(1)] 求$X$和$Y$的联合概率密度。
\item[(2)] 设含有$a$的二次方程为$a^{2}+2Xa+Y=0$，试求$a$有实根的概率。
\end{enumerate}

\setcounter{enumi}{20}
\item 设随机变量$(X,Y)$的概率密度为
\[
f(x,y)=
\begin{cases}
x+y, & 0<x<1,\ 0<y<1,\\[4pt]
0, & \text{其他}.
\end{cases}
\]
分别求(1) $Z=X+Y$、(2) $Z=XY$的概率密度。

\setcounter{enumi}{23}
\item 设随机变量$(X,Y)$的概率密度为
\[
f(x,y)=
\begin{cases}
\dfrac{1}{2}(x+y)\mathrm{e}^{-(x+y)}, & x>0,\ y>0,\\[6pt]
0, & \text{其他}.
\end{cases}
\]
\begin{enumerate}
\item[(1)] 问$X$和$Y$是否相互独立？
\item[(2)] 求$Z=X+Y$的概率密度。
\end{enumerate}

\setcounter{enumi}{28}
\item 设随机变量$(X,Y)$的概率密度为
\[
f(x,y)=
\begin{cases}
b\mathrm{e}^{-(x+y)}, & 0<x<1,\ 0<y<\infty,\\[4pt]
0, & \text{其他}.
\end{cases}
\]
\begin{enumerate}
\item[(1)] 试确定常数$b$。
\item[(2)] 求边缘概率密度$f_{X}(x),f_{Y}(y)$。
\item[(3)] 求函数$U=\max\{X,Y\}$的分布函数。
\end{enumerate}

\setcounter{enumi}{33}
\item 设$X,Y$是相互独立的随机变量，$X\sim \pi(\lambda_{1})$，$Y\sim \pi(\lambda_{2})$。证明
\[
Z=X+Y\sim \pi(\lambda_{1}+\lambda_{2}).
\]
\end{enumerate}

\end{document}